%!TEX root =  tfg.tex
\chapter{Seguridad}

\begin{quotation} [Pensador] {Confuncio (551 a.C. --479 a.C.)}
Cuando el objetivo te parezca difícil, no cambies de objetivo; busca un nuevo camino para llegar a él.
\end{quotation}

\begin{abstract}
En un mundo completamente digitalizado como en el que vivimos nuestros datos se vuelve un bien que muchos podrían querer vulnerar por lo que la seguridad se vuelve un elemento primordial en cualquier servicio.
\end{abstract}


\newpage
\section{Seguridad en Javascript}

Al producir código para una entidad como la Universidad se Sevilla uno de los principales objetivos que se deben perseguir es la seguridad de los datos de los usuarios. Es por ello que al hacer uso de un lenguaje de programación como Javascript que puede ser manipulado y se pueden introducir datos maliciosos en nuestra Base de Datos.

Tendremos que realizar un análisis previo para buscar qué podemos hacer como desarrolladores, qué debemos evitar y qué es inevitable y está fuera de nuestro contro.

Comencemos con lo que podemos hacer nosotros como desarrolladores:

\section{Buenas prácticas}

Para tratar de garantizar los seguieremos las siguientes prácticas:

\begin{itemize}

    \item Haremos uso de la herramienta AuditJS centrada en la búsqueda de vulnerabilidades. Una vez que se haya implementado una funcionalidad esta se lanzará en el código producido para garantizar que no existen vulneravilidades.
    \item A la hora de implementar formularios realizaremos validaciones en el backend tratando de descartar todo dato malicioso que nos llegue.
    \item Desinfectaremos y codificaremos las entradas del usuario para proteger nuestro servicio contra los ataques de inyección de consultas.
    \item Realizaremos análisis y documentaremos todas y cada una de las dependencias que implementemos para garantizar que no tienen vulnerabilidades conocidas.
    \item Implementaremos configuraciones como el atributo Http-Only en elementos como las Cookies para que el codigo JavaScript del lado del usuario no pueda robar el ID de sesión así como cualquier otra información confidencial del usuario.
    \item Haremos uso de tokens CSRF para evitar ataques de petición de sitios cruzados.    
    
\end{itemize}

Además de esto realizaremos otro tipo de configuraciones en nuestro servicio para evitar los ataques de DoS (Denegación de Servicio) como:

\begin{itemize}

    \item Limitaremos el número de conexiones que puede desarrollar un usuario.
    \item Limitaremos los lugares desde los cuales se puede conectar un usuario a nuestro servicio.

\end{itemize}
